The correlation of the current reference codes reads
$T_{\text{melt}} = 3113.15 - 5\,Bu/10000$, with $Bu$ in MWd/tU. It
\textbf{supersedes} the MATPRO one: the revision, by a factor of \num{6.4}, was
proposed by Popov \textit{et al.} (2000) on the basis of experimental data from
Adamson \textit{et al.} (1985) and Komatsu \textit{et al.}, the steep MATPRO
slope having proved too pessimistic~[11]. It is therefore the value of
$-0.5$ K/GWd/tU that has been adopted.

\vspace{0.3cm}
\noindent
\textit{Practical significance.} At 50 GWd/tU, the discrepancy between the two
correlations reaches \SI{135}{\kelvin} on the limit. On the reference case of
this report, where the centreline temperature peaks at \SI{1467}{\kelvin}, both
give the same verdict with a comfortable margin. The question only becomes
discriminating for a rod pushed to high burnup — that is, precisely the domain
where safety analysis has the most value.

%---------------------------------------------------------------- Appendix E --
\section{Training dataset of the surrogate}
\label{ann:dataset}

\noindent Full design of experiments: \num{5} values of linear heat rate
$\times$ \num{5} irradiation durations, that is \num{25} OFFBEAT simulations,
all of them carried through. Total cost: about \num{30} minutes of computation.

\vspace{0.3cm}
\noindent
\begin{tabularx}{\textwidth}{@{}Xll@{}}
\toprule
\textbf{Variable} & \textbf{Domain explored} & \textbf{Role}\\
\midrule
Linear heat rate & 12, 18, 24, 30, \SI{36}{\kilo\watt\per\metre} & input\\
Irradiation duration & \SIrange{7.9e6}{6.3e7}{\second} (3 months to 2 years)
  & input\\
\midrule
Maximum temperature & \SIrange{930}{2080}{\kelvin} & output\\
Maximum creep strain & \numrange{9.8e-4}{1.1e-2} & output\\
Minimum gap width & \SIrange{-13.2}{31.8}{\micro\metre} & output\\
\bottomrule
\end{tabularx}

\vspace{0.4cm}
\noindent
\textit{Choice of the durations.} The five values are placed just short of the
tabulated points of the template's power history. A duration that reaches or
exceeds the last time marker makes the solver fail, for want of a following
marker to drive the adaptive time step.

\vspace{0.3cm}
\noindent
\textit{Sign of the outputs.} A negative gap width signals an interference, that
is, an established pellet-cladding contact. It is precisely because the domain
covers this change of sign that the surrogate can learn to date the onset of
PCMI — which the first version, limited to durations of a few thousand seconds,
did not allow (§\ref{sec:validation-emulateur}).

\vspace{0.3cm}
\noindent
\textit{Range of validity.} The surrogate is reliable only within this domain.
Any prediction outside it is extrapolation; this is why the uncertainty of the
Gaussian process is systematically displayed next to the predicted value.

%---------------------------------------------------------------- Appendix F --
\section{Catalogue of faults in the self-healing benchmark}
\label{ann:bench}

\noindent Each fault is injected into a \emph{healthy} case generated
beforehand, so that the corruption is the only variable. The geometric faults
are injected after creation, deliberately to bypass input validation, the
purpose of the benchmark being to measure self-healing and not the upstream
safeguard.

\vspace{0.3cm}
\noindent
\begin{tabularx}{\textwidth}{@{}>{\raggedright\arraybackslash}p{3.5cm}Xcc@{}}
\toprule
\textbf{Category} & \textbf{Injected fault} &
\textbf{\begin{tabular}{@{}c@{}}Correction\\ in base\end{tabular}} &
\textbf{\begin{tabular}{@{}c@{}}Cause\\ treatable\end{tabular}}\\
\midrule
\multirow{2}{=}{Control} & Nominal case \SI{25}{\kilo\watt\per\metre}
  & --- & ---\\
 & Nominal case \SI{18}{\kilo\watt\per\metre} & --- & ---\\
\addlinespace
\multirow{3}{=}{Numerical} & Power \num{200} times nominal & yes & no\\
 & Power \num{100} times nominal & yes & no\\
 & Negative power & yes & no\\
\addlinespace
\multirow{3}{=}{Physical consistency} & Pellet wider than the inside of the
  cladding & no & no\\
 & Negative cladding thickness & no & no\\
 & Zero pellet radius & no & no\\
\addlinespace
\multirow{3}{=}{Missing file} & Main solver dictionary deleted
  & no & no\\
 & Initial temperature field deleted & no & no\\
 & Control dictionary deleted & no & no\\
\addlinespace
\multirow{3}{=}{Invalid entry} & Non-numeric simulated duration & no & no\\
 & Duration beyond the power history & \textbf{yes} & \textbf{yes}\\
 & Negative time step & no & no\\
\bottomrule
\end{tabularx}

\vspace{0.4cm}
\noindent
\textit{Reading the last two columns.} ``Correction in base'' indicates that the
knowledge base associates a remedy with the expected pattern; ``cause treatable''
indicates that this remedy can, by its nature, address the root cause. A single
fault in the catalogue meets both conditions, and it is the only one that
self-healing resolves (§\ref{sec:epreuve}).

%---------------------------------------------------------------- Appendix G --
\section{Worked example: a complete session}
\label{ann:session}

\noindent Transcript of a real session, captured after the corrections of
section~\ref{sec:angle-mort}. Tool results are abridged; the sequence of calls
is complete. The session was conducted in French; the request is translated
here.

\vspace{0.4cm}
\noindent\textbf{User request}
\begin{quote}
\itshape
Create a PWR fuel rod case in \texttt{/tmp/demo\_agent} with a linear heat rate
of 22 kW/m over about one year, run the simulation, then tell me whether the
safety margins are met.
\end{quote}

\vspace{0.2cm}
\noindent
\begin{tabularx}{\textwidth}{@{}>{\raggedright\arraybackslash}p{0.7cm}
                                >{\raggedright\arraybackslash}p{3.5cm}X@{}}
\toprule
\textbf{\#} & \textbf{Tool called} & \textbf{Result}\\
\midrule
1 & \texttt{input\_creator} & Case created, but the agent placed
  \texttt{template\_name} \emph{inside} the parameters: the key is reported as
  ignored\\
\addlinespace
2 & \texttt{offbeat\_executor} & \textcolor{cAlert}{Rejected}: ``is not a valid
  OpenFOAM case''. The agent had launched the execution \emph{in parallel} with
  the creation, before the case existed\\
\addlinespace
3 & \texttt{input\_creator} & The agent corrects itself: \texttt{template\_name}
  is moved back outside the parameters. Case created\\
\addlinespace
4 & \texttt{offbeat\_executor} & Mesh generated, simulation successful on the
  first attempt\\
\addlinespace
5 & \texttt{data\_processor} & Axial profile \SIrange{968.7}{1263.5}{\kelvin};
  radial profile centre \SI{1220.6}{\kelvin}, periphery \SI{601.6}{\kelvin};
  hoop stress \SI{-38.81}{\mega\pascal}; PCT \SI{634.5}{\kelvin};
  three figures generated\\
\bottomrule
\end{tabularx}

\vspace{0.4cm}
\noindent
\textbf{What this session shows.} Step~2 is the most instructive: the agent
makes a sequencing error — it launches the solver before the case exists — and
it is the \emph{deterministic safeguard} of the execution tool, not the language
model, that intercepts it. The agent reads the error message, infers the
correction and resumes. This is exactly the division of roles the architecture
aims at: the model orchestrates, the tools validate.

\vspace{0.3cm}
\noindent
\textbf{What it also shows.} The request comprised four steps; the agent carried
out three. It announced that it was going to launch the safety analysis, then
ended its turn without doing so. Long chaining therefore remains imperfect,
something the tool-selection benchmark did not measure either: it evaluates only
the first call, never the conduct of a sequence.

\end{document}
